\chapter{척도, 끈, 그리고 26차원}

\lectureoverview{끈 스펙트럼의 간격은 무엇으로 정해질까? 답은 장력이다. 장력의 차원은 질량의 제곱이다. 기본 끈에 적용하면 이 사실은 플랑크 단위, 직접 여기의 현실적인 불가능성, 강한 중력의 출현으로 이어진다. 잠시 쉬었다가 보손 열린끈의 바닥 상태로 돌아오자. 영점 에너지를 조절하면 역사적으로 24개의 횡방향과 그에 따른 26차원 시공간을 얻는 \emph{쉬운} 경로가 나온다.}

\section{빠뜨렸던 척도}

지금까지 우리는 끈의 진동과 회전을 이야기하면서 여기 한 단위가 무엇을 뜻하는지는 말하지 않았다. 하드론을 끈으로 그려 볼 수 있고 그 진동자들을 여기할 수는 있지만, 한 상태에서 다음 상태로 넘어가는 간격은 얼마나 될까? 그 간격은 무엇으로 정해질까?

기본적인 끈 모형에는 차원을 가진 매개변수가 하나뿐이다. 이를 찾으려면 늘어난 보통의 용수철부터 살펴보자. 양 끝점 사이의 거리가 \(\ell\)이면 훅의 법칙에 따라 다음을 얻는다.
\begin{equation}
  E_{\mathrm{spring}}=\frac{k\ell^2}{2}.
  \label{eq:l05-hooke}
\end{equation}
처음 작성한 뒤 수업 중에 계수 1/2이 추가되었다. 이 도입부의 추산에서는 이런 수치 계수에 본질적으로 의존하지 않는다.

\begin{figure}[H]\centering
  \includegraphics[width=0.72\textwidth]{lecture_05_figure_01.png}\caption{늘어난 구간과 훅의 법칙에 따른 퍼텐셜 에너지.}\lectureframe{00:02:28}\label{fig:l05-hooke-board}\end{figure}

\eqref{eq:l05-hooke}의 에너지는 비상대론적이다. 앞선 강의에서 사용한 무한 운동량 기술에서 내부의 비상대론적 에너지는 \(m\)이 아니라 \(m^2\)에 비례한다. 따라서 도식적으로 다음과 같이 쓴다.
\begin{equation}
  m^2\sim k\ell^2,
  \qquad
  m\sim \sqrt{k}\,\ell.
  \label{eq:l05-mass-length}
\end{equation}
끈의 정지계에서는 질량이 에너지와 같으므로, 두 번째 관계는 다음과 같은 꼴이다.
\begin{equation}
  E_{\mathrm{rest}}\sim T_s\ell,
  \qquad
  T_s\equiv\sqrt{k}.
  \label{eq:l05-tension}
\end{equation}
계수 \(T_s\)가 끈의 장력이다. 이는 단위 길이당 저장된 에너지다.

\begin{figure}[H]\centering
  \includegraphics[width=0.78\textwidth]{lecture_05_figure_02.png}\caption{판서에서 \(m^2\sim k\ell^2\)를 거쳐 길이에 비례하는 질량으로 넘어간다.}\lectureframe{00:05:42}\label{fig:l05-mass-length-board}\end{figure}

이 표기법은 강의의 용수철 비유에 맞춘 것이다. 표준 상대론적 끈 표기법에서는 보통 다음과 같이 쓴다.
\begin{equation}
  T_s=\frac{1}{2\pi\alpha'},
  \label{eq:l05-alpha-prime}
\end{equation}
따라서 끈의 레벨이 증가하면 \(m^2\)도 \(1/\alpha'\) 정도, 곧 장력 정도만큼 증가한다.
\editorialnote{강의에서 쓰는 유효한 \(k\)는 \(\sqrt{k}\)가 장력의 역할을 하도록 정해져 있다. 관례에 의존하지 않는 진술은 \(T_s\)의 차원이 질량의 제곱이며 레제 기울기가 \(\alpha'\propto T_s^{-1}\)를 만족한다는 것이다.}

\subsection{차원과 진동자 간격}

\(c=\hbar=1\)로 두자. 그러면 에너지의 차원은 역길이다.
\begin{equation}
  [E]=L^{-1}.
\end{equation}
장력은 단위 길이당 에너지이므로,
\begin{equation}
  [T_s]=L^{-2}=E^2.
  \label{eq:l05-tension-dimension}
\end{equation}
그래서 끈 스펙트럼의 간격은 질량 제곱으로 나타내는 것이 자연스럽다.

\begin{figure}[H]\centering
  \includegraphics[width=0.68\textwidth]{lecture_05_figure_03.png}\caption{자연 단위계의 판서: \(c=\hbar=1\)에서 에너지는 역길이이고, 진동자 진동수는 용수철의 척도로 정해진다.}\lectureframe{00:06:05}\label{fig:l05-natural-units-board}\end{figure}

보통의 진동자를 보면 같은 점을 다른 관점에서 확인할 수 있다. 질량이 \(\mu\)인 용수철에서는
\begin{equation}
  \omega=\sqrt{\frac{k}{\mu}},
  \qquad
  \Delta E=\hbar\omega.
\end{equation}
용수철이 더 뻣뻣할수록 진동수와 양자 간격이 커진다. 끈에는 이런 정상 모드가 무한히 많지만, 그 공통 척도는 \(T_s\)로 정해진다. 장력을 알면 끈 스펙트럼의 간격도 알 수 있다.

장력은 힘의 차원도 갖는다. 힘과 거리의 곱이 에너지이기 때문이다. 한쪽 끝을 고정하고 다른 쪽 끝에 추를 매달았다고 생각해 보자. \(E=T_s\ell\)인 이상적인 상대론적 끈은 길이와 무관하게 \(T_s\)로 정해지는 힘을 지탱할 수 있다. 관측된 여기 에너지로부터 추론한 하드론의 장력은, 일부러 엉뚱하게 비유하자면 트럭의 무게에 해당한다. 중간자에 트럭을 매달아 본 적은 없다. 측정된 스펙트럼에서 장력을 구하고, 그 장력으로 역학적 비유를 하는 것이다.

기본 끈은 훨씬 더 뻣뻣할 것이다. 전자, 광자, 중력자를 이루는 끈이라면 첫 진동자 여기 에너지가 관측된 입자 에너지보다 훨씬 높아야 한다. 같은 엉뚱한 중력 실험에 비유하면, 그 장력은 은하 질량에 해당하는 물체도 지탱할 수 있을 것이다. 이 비유의 요점은 공학적 가능성이 아니다. 엄청난 장력은 높은 여기 에너지와 매우 작은 바닥 상태 요동을 모두 뜻한다.

\section{보편 상수로 만드는 단위}

얼마나 큰 척도여야 이렇게 엄청날까? 자연스러운 기준은 플랑크 척도다. 값을 제시하는 대신 직접 구성해 보자.

미터, 킬로그램, 초는 실험실에서 쓰기에 훌륭한 단위지만, 그 기원은 인간의 규모에 맞춰져 있다. 미터는 천이나 밧줄을 재기에 편리하고, 피트는 신체 부위의 이름에서 왔으며, 초는 하루를 다루기 좋은 크기로 나눈 단위다. 양성자 반지름을 1로 두면 일부 핵물리학 공식은 간단해지겠지만, 양성자는 특정한 복합 입자 하나일 뿐이다. 그런 선택으로 물리학 전체가 간단해지지는 않는다.

차원 상수 세 개는 훨씬 더 보편적으로 쓰인다.
\begin{itemize}
  \item \(c\)은 모든 물질계의 한계 속도를 정한다.
  \item \(\hbar\)는 모든 양자계의 불확정성 관계에 나타난다.
  \item \(G\)는 중력을 모든 형태의 질량-에너지와 결합한다.
\end{itemize}
\(c=\hbar=G=1\)인 단위를 선택하는 일은 이 상수들로 길이, 시간, 질량 단위를 구성하는 것과 같다.

이 상수들의 보통 차원은 다음과 같다.
\begin{equation}
  \begin{aligned}
    [c]&=LT^{-1},
    & [\hbar]&=ML^2T^{-1},\\
    [G]&=L^3M^{-1}T^{-2}.
  \end{aligned}
  \label{eq:l05-dimensions-constants}
\end{equation}
가운데 관계는 \(\Delta x\,\Delta p\gtrsim\hbar\)에서 바로 나온다. 마지막 관계는 예를 들어 다음 식에서 얻을 수 있다.
\begin{equation}
  a=\frac{GM}{r^2}.
\end{equation}

\begin{figure}[H]\centering
  \begin{minipage}[t]{0.48\textwidth}\centering
    \includegraphics[width=\textwidth]{lecture_05_figure_04.png}\par\smallskip\footnotesize \(c\)\(\hbar\)의 단위 (a).
  \end{minipage}\hfill
  \begin{minipage}[t]{0.48\textwidth}\centering
    \includegraphics[width=\textwidth]{lecture_05_figure_05.png}\par\smallskip\footnotesize \(G\)의 단위: 뉴턴 가속도에서 구하기 (b).
  \end{minipage}\caption{수업에서 모은 차원 분석의 재료.}\lectureframe{00:25:43--00:26:40}\label{fig:l05-dimensional-analysis-board}\end{figure}

\subsection{먼저 플랑크 면적}

\(p,q,r\)의 거듭제곱을 찾아 다음 식이
\begin{equation}
  G^p\hbar^q c^r
\end{equation}
\(L^2\)의 차원을 갖게 하자. 먼저 면적을 구하면 초기에 제곱근을 취하지 않아도 되고, 나중에 중력에서 면적이 특별한 역할을 하리라는 점도 미리 드러난다. \eqref{eq:l05-dimensions-constants}에서 질량과 시간의 차원을 소거하면
\begin{equation}
  -p+q=0,
  \qquad
  -2p-q-r=0.
\end{equation}
남은 길이 차원은 다음 조건을 만족해야 한다.
\begin{equation}
  3p+2q+r=2.
\end{equation}
따라서 \(q=p\), \(r=-3p\), \(p=1\)이다. 플랑크 면적은
\begin{equation}
  \ell_P^2=\frac{G\hbar}{c^3},
  \label{eq:l05-planck-area}
\end{equation}
그러므로
\begin{equation}
  \ell_P=\sqrt{\frac{G\hbar}{c^3}},
  \qquad
  t_P=\frac{\ell_P}{c},
  \qquad
  m_P=\sqrt{\frac{\hbar c}{G}}.
  \label{eq:l05-planck-units}
\end{equation}

\begin{figure}[H]\centering
  \includegraphics[width=0.78\textwidth]{lecture_05_figure_06.png}\caption{플랑크 면적과 길이에 대한 차원 분석의 완성된 논증.}\lectureframe{00:31:34}\label{fig:l05-planck-area-board}\end{figure}

차수 추산의 정확도에서는
\begin{equation}
  \begin{aligned}
    \ell_P&\sim10^{-35}\,\mathrm{m},
    &t_P&\sim10^{-43}\,\mathrm{s},\\
    m_P&\sim10^{-8}\,\mathrm{kg},
    &m_Pc^2&\sim10^{19}\,\mathrm{GeV}.
  \end{aligned}
  \label{eq:l05-planck-numbers}
\end{equation}

\begin{figure}[H]\centering
  \includegraphics[width=0.72\textwidth]{lecture_05_figure_07.png}\caption{차수 추산으로 기록한 세 플랑크 척도.}\lectureframe{00:35:12}\label{fig:l05-planck-values-board}\end{figure}

길이와 시간은 상상할 수 없을 만큼 작지만, 질량은 입자의 기준에서 놀라울 정도로 거시적이며 고운 먼지 한 톨의 질량과 비슷하다. 정지 에너지는 수 기가줄 정도로, 휘발유 한 탱크에 든 화학 에너지와 비슷하다. 반면 관측 가능한 우주의 크기는 대략 \(10^{60}\) 플랑크 길이이고, 그 나이는 대략 \(10^{60}\) 플랑크 시간이다. 이런 비교는 단위에 물리적 감각을 부여하기 위한 것이지 수치의 정밀도를 높이기 위한 것이 아니다.

특성 끈 길이 \(\ell_s=\sqrt{\alpha'}\)가 \(\ell_P\)에 가깝다면, 끈의 양자적 크기는 \(\ell_P\)에, 진동 주기는 \(t_P\)에, 첫 번째 질량 있는 여기는 \(m_P\)에 가까울 것이다. 강의에서는 그 값에 10배, 100배 또는 그 이상의 차이가 있을 수 있다고 본다. 이는 차원 항등식이 아니라 예상이다.

\begin{classroomqa}\audiencequestion{플랑크 척도 질량은 입자의 바닥 상태에서 증가한 양인가? \lecturetimestamp{00:38:37}}\lecturerresponse{간격이 일정한 양은 질량 제곱이다. 바닥 상태가 가벼우면 첫 번째 끈 여기는 \(m^2\)가 \(m_P^2\) 정도만큼 더 크므로, 그 질량과 필요한 질량중심 에너지는 \(m_P\) 정도다.}\end{classroomqa}

\section{두 척도는 얼마나 가까워야 할까?}

\begin{classroomqa}\audiencequestion{차원 분석으로는 \(\ell_P\)를 얻는데, 왜 이것이 끈의 길이여야 하는가? \lecturetimestamp{00:40:21}}\lecturerresponse{그것은 차원 분석만의 결과가 아니라 중력과 끈의 구조가 관측되지 않았다는 사실에 근거한 물리적 추측이다. 끈 길이는 플랑크 길이보다 클 수 있으며, 그러면 에너지 척도는 더 낮다. 강의의 주장은 그 길이가 보통의 입자물리학적 길이는 아니어야 하고, 플랑크 길이보다 훨씬 작게 만들면 중력 붕괴에 부딪힌다는 것이다.}\end{classroomqa}

두 척도의 관계는 실제로 결합상수에 따라 달라진다. 완전한 콤팩트화에서는 \(\ell_P\), \(\ell_s\), 끈 결합상수, 콤팩트화 부피가 함께 들어간다. 보편적으로 성립하는 방정식 \(\ell_s=\ell_P\)은 없다. 강의에서 유효한 것은 척도에 관한 논증이다. 끈 효과는 매우 높은 에너지에서 나타나야 하며, 중력을 무시한 채 플랑크 길이보다 작은 영역에 국소화하는 문제를 논할 수는 없다.
\editorialnote{현대 끈 이론 구성에서 끈 길이가 4차원 플랑크 길이와 정확히 같을 필요는 없다. 둘 사이의 관계는 끈 결합상수와 콤팩트화 자료에 따라 달라진다.}

\begin{classroomqa}\audiencequestion{플랑크 길이를 유용한 최소 길이로 본다면, 플랑크 질량도 가능한 최소 질량인가? \lecturetimestamp{00:43:27}}\lecturerresponse{아니다. 보통 입자의 질량은 훨씬 작다. 플랑크 질량은 작은 천체의 중력 반지름이 양자 파장과 비슷해지는 척도다. 강의의 직관적 표현을 따르면 수명은 플랑크 시간 정도인 가장 가벼운 준고전적 블랙홀의 척도다. 마지막 진술을 완전한 양자 물체에 대한 정밀한 설명으로 오해해서는 안 된다.}\end{classroomqa}

\subsection{달리는 결합상수로 알 수 있는 것}

통상적인 양자장론이 직접 가속기 실험이 가능한 범위를 훨씬 넘어선 척도까지 유효하다는 간접 증거가 있다. 표준모형의 세 게이지 결합상수는 접근 가능한 에너지에서 측정되며, 그 에너지 의존성을 외삽할 수 있다. 강의의 개략적인 그래프에서는 \(U(1)\), \(SU(2)\), \(SU(3)\) 결합상수가 높은 통일 척도에서 서로 가까워진다. 그 척도는 플랑크 에너지보다 몇 자릿수 정도 낮을 수도 있다.

\begin{figure}[H]\centering
  \includegraphics[width=0.70\textwidth]{lecture_05_figure_08.png}\caption{개략적인 결합상수의 에너지 의존성 그래프. 세 게이지 곡선은 중력의 세기가 증가해 1 정도가 되기 전에 서로 가까워진다.}\lectureframe{00:47:18}\label{fig:l05-running-couplings-board}\end{figure}

이는 교차점에서 직접 실험한 결과가 아니라 외삽이다. 실험실 에너지와 통일 에너지 사이에 추가 입자나 문턱이 있으면 기울기가 달라질 수 있다. 신중하게 말할 수 있는 결론은 알려진 장론에서 보통 거리 규모의 끈 크기 하부 구조가 드러날 조짐이 없으며, 그럴듯한 끈 척도가 추론된 통일 척도와 플랑크 척도 사이의 엄청나게 넓은 구간 어딘가에 있다는 것뿐이다.

\begin{classroomqa}\audiencequestion{어떤 세 결합상수를 외삽하는가? \lecturetimestamp{00:48:28}}\lecturerresponse{그것들은 \(U(1)\), \(SU(2)\), \(SU(3)\)에 해당하는 게이지 결합상수, 즉 약전하, 약한 상호작용, 강한 상호작용의 결합상수다. 이 분류에서 전자기력은 독립된 네 번째 게이지 인자가 아니라 전약 혼합에서 나온다.}\end{classroomqa}

\begin{classroomqa}\audiencequestion{중력상수 \(G\) 자체가 더 높은 에너지를 탐침할수록 커지는가? \lecturetimestamp{00:49:30}}\lecturerresponse{\(G\)은 차원을 가진 양이므로, 그 자체로는 무차원 게이지 결합상수와 비교하기에 적절한 양이 아니다. 특성 에너지 \(M\)에서 무차원 중력 세기는 도식적으로 \(GM^2/(\hbar c)\)처럼 거동한다. 이 조합은 에너지와 함께 커져 플랑크 척도에서 1 정도가 된다.}\end{classroomqa}

힘의 법칙을 보면 이 비교가 분명해진다.
\begin{equation}
  F_{\mathrm{em}}\sim\frac{e^2}{r^2},
  \qquad
  F_{\mathrm{grav}}\sim\frac{GM^2}{r^2}.
  \label{eq:l05-force-laws}
\end{equation}
따라서 자연 단위계에서는 \(e^2\)을 \(G\)과 비교하는 것이 아니라 \(GM^2\)과 비교한다.

\begin{figure}[H]\centering
  \includegraphics[width=0.74\textwidth]{lecture_05_figure_09.png}\caption{무차원 중력 세기를 확인하는 데 사용한 전자기력과 중력의 힘의 법칙.}\lectureframe{00:50:50}\label{fig:l05-force-laws-board}\end{figure}

\section{직접 끈 분광학은 왜 요원한가}

전자가 끈이라는 사실을 실험으로 어떻게 확인할 수 있을까? 하드론 분광학에서 한 것처럼 내부 모드를 여기할 만큼 세게 충돌시킨 다음, 질량이 끈의 패턴을 따르는 여러 상태의 탑을 찾아보면 된다. 첫 번째 질량 있는 레벨이 \(10^{19}\,\mathrm{GeV}\) 근처라면 직접 실험은 어떤 가속기로도 도달할 수 없는 범위에 있다.

장벽은 총에너지만이 아니다. 에너지를 아주 작은 충돌 매개변수 안에 집중해야 한다. 여기서 블랙홀이 등장한다.

\begin{classroomqa}\audiencequestion{이 에너지에서는 가속기가 작은 블랙홀을 만들기 시작하는가? \lecturetimestamp{00:50:57}}\lecturerresponse{질량중심 에너지를 중력 반지름 안에 집중하면 그렇게 되리라는 것이 직관적 예상이다. 항성 블랙홀은 중력 퍼텐셜 에너지가 무거운 별을 압축해 만들어진다. 미시적 블랙홀은 운동 에너지를 작은 충돌 영역에 집중시켜 만들어질 것이다. 4차원 중력에서는 필요한 척도가 플랑크 척도다.}\end{classroomqa}

규모를 가늠해 보자. 선형 가속기의 에너지 증가는 대략 길이에 비례한다. 강의에서는 대략 \(100\,\mathrm{GeV}\)를 SLAC 규모의 간단한 기준으로 삼는다. 같은 가속 기울기에서 \(10^{19}\,\mathrm{GeV}\)에 도달하려면 길이가 대략 \(10^{17}\)배가 되어 은하 규모의 장치가 필요하다.

그런 장치도 광도가 필요하다. 빔을 한 쌍씩 플랑크 길이 정밀도로 조준하는 것이 아니다. 엄청난 수의 입자를 충돌시켜 그중 몇 개가 원하는 충돌 매개변수에서 만나게 한다. 플랑크 에너지에 엄청난 충돌률을 곱하면 터무니없을 정도로 큰 전력이 필요하다. 이 추산은 일부러 거칠게 한 것이지만 결론은 분명하다. 플랑크 척도에서 직접 끈 분광학을 수행할 실험을 현실적으로 설계할 수는 없다.

간접 증거는 더 구조적이어야 한다. 관측된 저에너지 입자와 상호작용을 구체적으로 예측해 검증할 수 있는 이론이거나, 끈 척도보다 훨씬 낮은 에너지까지 남아 있는 콤팩트화 효과여야 한다. 강의는 그런 결정적인 구성이 이미 마련되어 있다고 주장하지 않는다.

\subsection{추가 차원과 국소화 한계}

\begin{classroomqa}\audiencequestion{추가 차원의 크기도 플랑크 길이 정도여야 하는가? \lecturetimestamp{00:56:23}}\lecturerresponse{이 기본적인 그림에서는 추가 차원이 극도로 작을 것으로 예상하지만, 이 논증만으로 크기가 정확히 정해지지는 않는다. 콤팩트 차원이 낮은 에너지에서 탐사할 만큼 크다면, 새로운 칼루차--클라인 모드와 고차원적 거동이 이미 우리가 관측하는 장론을 바꾸었을 것이다. 따라서 장론이 성공적으로 외삽된다는 점은 그런 차원의 크기에 상한을 둔다.}\end{classroomqa}

\begin{classroomqa}\audiencequestion{끈이나 콤팩트 차원의 크기가 \(\ell_P\)보다 여러 자릿수나 더 작을 수 없는 이유는 무엇인가? \lecturetimestamp{00:57:48}}\lecturerresponse{더 짧은 거리를 분해하려면 더 큰 운동량이 필요하다. 플랑크 척도를 넘어서면 국소화에 필요한 에너지의 중력 반지름이 탐사 영역과 비슷해지고, 그보다 더 커진다. 더 작은 보통의 끈을 분해하기도 전에 실험이 블랙홀을 만들어 낸다.}\end{classroomqa}

불확정성 원리로 추산할 수 있다. \(\Delta x\)까지 국소화하려면

\begin{classroomqa}\audiencequestion{플랑크 길이와 플랑크 시간을 정한 뒤 플랑크 질량은 물리적으로 어떻게 정해지는가? \lecturetimestamp{00:59:04}}\lecturerresponse{길이 척도를 에너지로 바꾸려면 \(\hbar\) 또는 \(G\)를 사용하면 된다. 불확정성 원리에 따르면 두 입자를 \(\ell_P\) 안에 국소화하려면 운동량이 \(\hbar/\ell_P\) 정도여야 하며, 따라서 에너지는 \(m_Pc^2\) 정도다. 바꾸어 말하면, 양자 길이와 중력 길이가 비슷해지는 질량이다.}\end{classroomqa}

\begin{equation}
  \Delta p\gtrsim\frac{\hbar}{\Delta x},
  \qquad
  E\sim c\,\Delta p.
  \label{eq:l05-localization}
\end{equation}
\(\Delta x=\ell_P\)일 때 이 에너지는
\begin{equation}
  E\sim\frac{\hbar c}{\ell_P}=m_Pc^2.
\end{equation}
동시에 슈바르츠실트 반지름
\begin{equation}
  r_s\sim\frac{GE}{c^4}
\end{equation}
도 수치 계수를 제외하면 \(\ell_P\) 정도다. 양자 국소화와 중력 붕괴가 같은 척도에서 만난다.

\begin{figure}[H]\centering
  \includegraphics[width=0.74\textwidth]{lecture_05_figure_10.png}\caption{국소화 논증에서 다시 살펴본 플랑크 길이와 질량의 공식.}\lectureframe{01:00:25}\label{fig:l05-localization-board}\end{figure}

이로써 척도 설정 논증을 마친다.

\section{바닥 상태는 광자보다 낮아야 한다}

휴식 뒤에는 두 가지 길이 있다. 등각 변환을 시작하거나, 임계 차원으로 가는 더 오래되고 쉬운 경로를 제시하는 것이다. 후자를 택하되 주의할 점이 있다. 이 논증은 쉽지 않고 가장 근본적인 이유를 설명하지도 않는다. 그렇지만 옳은 논증이며 역사적 발견에 가까운 경로다.

보손 열린끈으로 돌아가자. 횡방향 생성 연산자를 한 번 작용시키면 편극 상태 두 개가 생긴다. 세 번째 종방향 상태가 없으므로 이 여기는 질량 없는 벡터여야 한다. 한 모드-\(1\) 양자가 \(m^2\)를 1만큼 높이는 단위계에서는
\begin{equation}
  m_{\mathrm{ground}}^2=-1,
  \qquad
  m_{\mathrm{first}}^2=-1+1=0.
  \label{eq:l05-open-intercept}
\end{equation}
타키온 바닥 상태는 보손 이론에 여전히 남아 있는 문제다. 여기서는 이를 인정하고 잠시 제쳐 둔다. 당장 묻는 것은 \eqref{eq:l05-open-intercept}의 음의 단위가 어디서 오는가이다.

끈의 각 모드는 조화 진동자다. \(\hbar=1\)일 때 진공 에너지는
\begin{equation}
  E_{0,n}=\frac{\omega_n}{2}.
\end{equation}
끈에서는 \(\omega_n=n\)이므로 횡방향 좌표 하나가 형식적으로 기여하는 양은
\begin{equation}
  E_0^{(1)}
  =\frac12\sum_{n=1}^{\infty}n.
  \label{eq:l05-formal-zero-point}
\end{equation}

\begin{figure}[H]\centering
  \includegraphics[width=0.70\textwidth]{lecture_05_figure_11.png}\caption{모드 진동수를 되짚은 뒤 적은 영점 합 \(\sum n/2\).}\lectureframe{01:06:18}\label{fig:l05-zero-point-board}\end{figure}

이제 난점이 분명해진다. \eqref{eq:l05-formal-zero-point}의 모든 항은 양수이고 합은 발산한다. 그런데 어떻게 유한한 음의 절편을 줄 수 있을까?

\section{스펙트럼을 읽기 전에 운동을 빼기}

합을 조절하기 전에 내부 에너지를 질량 제곱과 동일시한 이유를 떠올려 보자. 큰 종방향 운동량 \(P>0\)을 가진 입자에 대해,
\begin{equation}
  E=\sqrt{P^2+m^2}
   =P+\frac{m^2}{2P}
   +O(P^{-3}).
  \label{eq:l05-large-p-energy}
\end{equation}

\begin{figure}[H]\centering
  \includegraphics[width=0.70\textwidth]{lecture_05_figure_12.png}\caption{전체 운동량과 내부 질량 항을 분리하는 큰 \(P\) 전개.}\lectureframe{01:07:35}\label{fig:l05-fast-energy-board}\end{figure}

첫 번째 항은 \(z\)축을 따라 움직이는, 관심 대상이 아닌 전체 운동을 나타낸다. 이 항은 보존되며 \(P\)가 같은 모든 내부 상태에 공통이다. 이 항을 빼고 재조정하면 운동학적 시간 지연이 사라진다.
\begin{equation}
  \begin{aligned}
    E-P&=\frac{m^2}{2P}+O(P^{-3}),\\
    2P(E-P)&=m^2+O(P^{-2}).
  \end{aligned}
  \label{eq:l05-internal-energy}
\end{equation}

원자가 물리적 그림을 제공한다. 수소 원자를 부스트하면 총에너지는 커지지만, 원자의 시계가 느려지므로 실험실 시간으로 잰 두 내부 준위 사이의 에너지 차이는 줄어든다. \eqref{eq:l05-large-p-energy}에서 계수 \(1/P\)가 바로 그 억제를 나타낸다. \(2P\)를 곱하면 불변 질량 제곱 차이가 복원된다.

따라서 바닥 상태의 날것 그대로의 에너지가 반드시 문자 그대로 \(-1\)일 필요는 없다. 전체 종방향 에너지와 관련된 큰 상태 독립 항에 유한한 내부 기여가 더해져 있을 수 있다. 강의에서는 진동자 합으로 돌아가기 전에 부스트된 관계식을 다시 써 두 항을 드러낸다.

\begin{figure}[H]\centering
  \includegraphics[width=0.76\textwidth]{lecture_05_figure_13.png}\caption{큰 상태 독립 종방향 기여와 나란히 나타낸 불변 질량 항.}\lectureframe{01:11:34}\label{fig:l05-infinite-piece-board}\end{figure}

양자장론에서는 무한한 진공 에너지가 낯설지 않다. 발산하는 상태 독립 상수와 관측 가능한 에너지 차이를 분리하는 경우가 많다. 강의는 이를 방 안의 전자기장의 영점 에너지와 비교하며, 이 절차가 만족스럽지는 않지만 표준적이라고 솔직히 말한다. 여기서는 같은 직관적 방법을 따라 발산 부분을 분리하고 명확한 유한 잔여항이 남는지 살펴본다.

\section{합 조절하기}

양의 조절 변수 \(\epsilon\)를 도입하자.
\begin{equation}
  S_1(\epsilon)
  =\sum_{n=1}^{\infty}n e^{-n\epsilon}.
  \label{eq:l05-regulated-sum}
\end{equation}
모든 \(\epsilon>0\)에 대해 지수 함수가 결국 인자 \(n\)를 압도하므로 합은 수렴한다. \(\epsilon\)를 0에 가깝게 보내면 절단이 제거된다.

\begin{figure}[H]\centering
  \includegraphics[width=0.78\textwidth]{lecture_05_figure_14.png}\caption{\(1+2+3+\cdots\)에 항별로 지수 조절 인자를 적용한 결과.}\lectureframe{01:16:42}\label{fig:l05-regulator-board}\end{figure}

미분을 이용하면 인자 \(n\)를 만들 수 있다.
\begin{equation}
  S_1(\epsilon)
  =
  -\frac{d}{d\epsilon}
  \sum_{n=1}^{\infty}e^{-n\epsilon}.
  \label{eq:l05-derivative-trick}
\end{equation}

\begin{figure}[H]\centering
  \includegraphics[width=0.76\textwidth]{lecture_05_figure_15.png}\caption{절단 변수에 대해 미분하면 \(-n\)가 따라 나온다.}\lectureframe{01:18:05}\label{fig:l05-derivative-board}\end{figure}

남은 급수는 기하급수다.
\begin{equation}
  S_0(\epsilon)
  \equiv
  \sum_{n=1}^{\infty}e^{-n\epsilon}
  =
  \frac{e^{-\epsilon}}{1-e^{-\epsilon}}
  =
  \frac{1}{e^\epsilon-1}.
  \label{eq:l05-geometric-sum}
\end{equation}

\begin{figure}[H]\centering
  \includegraphics[width=0.68\textwidth]{lecture_05_figure_16.png}\caption{닫힌 형태로 쓴 조절된 기하급수.}\lectureframe{01:19:12}\label{fig:l05-geometric-board}\end{figure}

작은 \(\epsilon\)에서 전개하자.
\begin{equation}
  S_0(\epsilon)
  =
  \frac{1}{\epsilon}
  -\frac12
  +\frac{\epsilon}{12}
  -\frac{\epsilon^3}{720}
  +O(\epsilon^5).
  \label{eq:l05-bernoulli-expansion}
\end{equation}
판서에서는 분자와 분모를 따로 전개하고, \(\epsilon\)의 한 거듭제곱을 묶은 뒤 역수를 전개해 이 결과를 얻는다. 이 계산에서는 수업 중 부호를 여러 차례 확인했으며, 간결한 식 \eqref{eq:l05-bernoulli-expansion}는 정리된 결과다.

\begin{figure}[H]\centering
  \includegraphics[width=0.78\textwidth]{lecture_05_figure_17.png}\caption{역수를 간단히 하기 전의 분자와 분모 전개.}\lectureframe{01:23:10}\label{fig:l05-expansion-board}\end{figure}

\(-d/d\epsilon\)를 적용하면 다음을 얻는다.
\begin{equation}
  S_1(\epsilon)
  =
  \frac{1}{\epsilon^2}
  -\frac{1}{12}
  +\frac{\epsilon^2}{240}
  +O(\epsilon^4).
  \label{eq:l05-regulated-result}
\end{equation}

\begin{figure}[H]\centering
  \includegraphics[width=0.78\textwidth]{lecture_05_figure_18.png}\caption{조절 계산의 끝에 얻은 발산 항 \(1/\epsilon^2\)과 유한한 잔여항 \(-1/12\).}\lectureframe{01:29:48}\label{fig:l05-minus-twelve-board}\end{figure}

이 결과는 농담 삼아 \(1+2+3+\cdots=-1/12\)이라고 요약하곤 한다. 이는 보통 의미의 수렴이 아니다. \eqref{eq:l05-regulated-result}는 이 특정 조절법에서 발산 항과 조절 변수에 무관한 유한 부분이 나온다는 뜻이다.
\begin{equation}
  \underbrace{\frac{1}{\epsilon^2}}_{\text{divergent}}
  \;+\;
  \underbrace{\left(-\frac{1}{12}\right)}_{\text{finite}}
  \;+\;O(\epsilon^2).
\end{equation}
강의에서는 첫 항을 이미 존재하는 상태 독립 부분에 흡수하고 두 번째 항을 남긴다. 또한 무한대라고 해서 저절로 0이 되는 것은 아니라는 파울리의 경고도 유지한다. 일관된 이론이라면 그 뺄셈이 허용되는 이유를 설명해야 한다.

\eqref{eq:l05-formal-zero-point}에서 1/2 계수를 복원하면 횡방향 좌표 하나의 기여는 다음과 같다.
\begin{equation}
  E_0^{(1)}=-\frac{1}{24}.
  \label{eq:l05-one-transverse-zero-point}
\end{equation}

\begin{figure}[H]\centering
  \includegraphics[width=0.68\textwidth]{lecture_05_figure_19.png}\caption{1/2 계수를 적용하면 유한한 \(-1/12\)은 횡방향 좌표당 \(-1/24\)이 된다.}\lectureframe{01:33:42}\label{fig:l05-minus-twenty-four-board}\end{figure}

\section{횡방향 수 세기}

횡방향 끈 좌표의 수를 \(D_\perp\)라 하자. 이들이 진공에 기여하는 전체 양은
\begin{equation}
  E_0^{\mathrm{open}}
  =-\frac{D_\perp}{24}.
  \label{eq:l05-open-vacuum}
\end{equation}
첫 번째 열린끈 진동자가 바닥 상태를 \(-1\)에서 질량 없는 레벨까지 올리려면 다음이 필요하다.
\begin{equation}
  -\frac{D_\perp}{24}=-1.
\end{equation}
따라서
\begin{equation}
  D_\perp=24.
\end{equation}
이 방향들은 큰 종방향 운동량에 수직이다. 여기에 종방향 공간 방향 하나와 시간을 더하면
\begin{equation}
  D=D_\perp+2=26.
  \label{eq:l05-critical-dimension}
\end{equation}

\begin{classroomqa}\audiencequestion{이 논증은 언제 끈 이론에 들어왔는가? \lecturetimestamp{01:35:27}}\lecturerresponse{이는 하드론 끈 이론 시기, 1969년경과 초기 1970년대에 등장했다. 레제 궤적은 이미 실험으로 알려져 있었다. 진동자 모형 자체를 탐구하면서 첫 번째 들뜬 상태에는 벡터 편극이 세 개가 아니라 두 개 있다는 사실이 드러났고, 이에 따라 그 상태는 질량이 없어야 했다. 필요한 절편으로부터 24개의 횡방향 진동자가 곧바로 도출되었다.}\end{classroomqa}

당시에는 믿기 어려울 만큼 이상하게 보였다. 조절된 진공 합이 어떻게 시공간의 차원을 알아낼 수 있을까? 그런데 여러 요소가 맞아떨어졌다. 질량 없는 벡터, 절편, 횡방향 진동자의 수가 모두 같은 값을 선택했다.

\subsection{닫힌끈으로 수를 확인하다}

닫힌끈이 첫 물리적 질량 없는 레벨에 도달하려면 좌이동 모드와 우이동 모드가 하나씩 필요하다. 따라서 그 바닥 상태는 0보다 2단위 아래에서 시작해야 한다.

\begin{classroomqa}\audiencequestion{횡방향이 24개라면 \(-1\)만큼 낮아질 뿐인데, 닫힌끈은 횡방향이 48개인 시공간에 살아야 하는가? \lecturetimestamp{01:37:47}}\lecturerresponse{아니다. 열린끈과 닫힌끈은 같은 시공간에 존재한다. 닫힌끈에는 이미 진동자 내용이 두 배로 들어 있다. 세계면의 각 방향으로 하나씩 탑이 돈다. 각 탑은 \(-24/24=-1\)를 기여하므로 합하면 \(-2\)이 된다.}\end{classroomqa}

따라서
\begin{equation}
  E_0^{\mathrm{closed}}
  =
  2\left(-\frac{D_\perp}{24}\right)
  =-2
  \qquad (D_\perp=24).
\end{equation}
각 부문에서 진동자 하나씩을 더하면 상태가 2단위 올라 질량 제곱이 0인 곳에 도달한다. 이는 앞선 강의에서 중력자 다중항이 나타난 바로 그 지점이다.

이 일치는 논증을 더 그럴듯하게 만들지만, 여전히 근본적인 설명은 아니다. 현대적인 양자화에서 정규순서화 상수, 로렌츠 대수 또는 비라소로 대수의 닫힘, 그리고 등각 이상 현상의 상쇄는 하나의 일관된 구조에 속한다. 이 구조가 절편과 임계 차원을 모두 결정하며, 유한 절단값의 나머지 항은 그 구조에 이르는 역사적인 출발점이다.
\editorialnote{조절자 계산은 정규순서화 상수를 휴리스틱하게 이해하는 관점이다. \(D=26\)을 더 근본적으로 유도하려면 등각 이상 현상의 상쇄 및 제약 대수의 일관성과 동등한 세계면 등각 불변성이 필요하다.}

이제 그보다 근본적인 접근으로 넘어가자. 시작하는 데 필요한 복소해석학은 아주 기본적인 내용뿐이다. 복소 좌표, 해석 함수, 그리고 코시--리만 방정식이 그것이다. 약속한 물리적 이미지는 터키식 태피---한쪽으로 당기고 다른 쪽으로 늘려, 중요한 구조를 무너뜨리지 않은 채 모양을 다시 만드는 것이다. 그 이미지의 바탕이 되는 수학이 등각 불변성이다.